\chapter{Semantic Analysis and IR\_TREE}
\label{chap:semantic}

Semantic analysis is the stage where the compiler decides what the program
means. In \epl{}, semantic analysis is implemented primarily by lowering the AST
into \irtree{}. The resulting representation is still expression-oriented, but it
is no longer source syntax. Every expression has a type, every local variable and
function has an ID, struct and array layouts are known, and several source
constructs have been rewritten into simpler core forms.

\section{Purpose of IR\_TREE}

\irtree{} is the main semantic representation of the compiler. It exists between
the AST and the lower CFG IR. It is needed because the AST is too source-like for
code generation, while the lower CFG IR is too low-level for convenient semantic
checking and compile-time evaluation.

The key properties of \irtree{} are:

\begin{itemize}
  \item each expression carries a static type;
  \item local variables are represented by \texttt{VariableId}, not by source
        strings;
  \item functions are represented by \texttt{FunctionId};
  \item loops are represented by \texttt{LoopId}, which connects
        \texttt{break} and \texttt{continue} to the correct loop;
  \item places are represented separately from values;
  \item arrays, structs, field access, and pointer dereference have explicit
        typed forms;
  \item source conveniences such as \texttt{while}, \texttt{for}, \texttt{\&\&},
        and \texttt{||} are lowered to simpler expression forms.
\end{itemize}

Because \irtree{} is readable, it is also used by the test harness as a snapshot
format. The snapshots are not a formal stable API, but they are very useful for
regression testing and for explaining the compiler in this thesis.

\section{Module Construction}

\texttt{Module::from\_ast} constructs a typed module from the AST. It performs
several passes:

\begin{enumerate}
  \item create a new type system context and seed the global type namespace with
        built-in types;
  \item collect struct definitions and compute their layouts;
  \item collect function declarations, including signatures, return types,
        variadic flags, and \texttt{@pure} annotations;
  \item lower each function body using the completed type and function
        namespaces;
  \item run semantic checkers;
  \item run basic tree optimizations;
  \item evaluate every remaining \texttt{comptime} expression and replace it
        with a constant.
\end{enumerate}

This ordering is important. Function declarations must be known before function
bodies are lowered so a function can call another function that appears later in
the file. Struct definitions are collected before function bodies, but their
fields are processed in source order. This means a struct definition can refer
only to types already known at that point.

\section{Type Resolution}

AST types are resolved by looking up identifier names in the type namespace.
Built-in types are inserted before source items are processed. Struct
definitions add new entries to the namespace. Pointer and array types are
constructed recursively.

Struct layout follows a simple alignment algorithm. For each field, the current
size is rounded up to the field's alignment, the field offset is recorded, and
the field size is added. The final struct size is rounded up to the maximum
field alignment. Arrays use the element layout and multiply by the length. Unit
and never have zero size. Booleans and \texttt{i8}/\texttt{u8} have one-byte
layout. \texttt{i32}/\texttt{u32} have four-byte layout, and
\texttt{i64}/\texttt{u64} have eight-byte layout.

\section{Name Resolution and Scope}

Function body lowering uses a \texttt{FunctionLoweringCtx}. The context stores
the function declaration, function namespace, function table, type system, type
namespace, and current lexical scope. A scope contains a map from source variable
names to \texttt{VariableId} and type pairs. It also stores optional loop
context for resolving \texttt{break} and \texttt{continue}.

When entering a block, the compiler pushes a new scope. When leaving the block,
it pops the scope. Lookup searches the current scope and then recursively
searches parent scopes. This implements lexical scoping and shadowing.

Function parameters are lowered into local variables at the beginning of the
function body. Each parameter receives a variable ID and a declaration in the
outer function block. The body begins with stores from \texttt{Argument(n)}
expressions into those locals. This design makes parameters behave like mutable
locals and avoids a second storage model for arguments.

\section{Expected Types and Local Inference}

\epl{} uses local, context-driven type inference. The compiler does not infer
function signatures, but it can infer the type of a local variable from its
initializer, infer unsuffixed integer literal types from the expected type, and
propagate expected types into several expression forms.

Expected types are passed into lowering for:

\begin{itemize}
  \item typed \texttt{let} initializers;
  \item assignment right-hand sides;
  \item fixed function call arguments;
  \item function return expressions;
  \item block final expressions;
  \item \texttt{if} branches;
  \item unnamed struct initializers;
  \item array initializers;
  \item \texttt{undefined};
  \item unsuffixed integer literals.
\end{itemize}

This approach gives useful inference while keeping the implementation simple.
For example, \texttt{let x: u64 = 0;} works because the expected type
\texttt{u64} flows into the literal \texttt{0}. Without that context, the same
literal would default to \texttt{i32}.

\section{Places and Values}

The compiler distinguishes places from values. A place is something that can be
assigned to or addressed: a local variable, a field of a place, an array element
of a place, or a dereferenced pointer. A value is the result of evaluating an
expression. Source syntax does not always expose this distinction, so semantic
lowering derives it.

For example, parsing \texttt{x.y} produces a field-access AST expression. During
semantic lowering it can become either a value field access or, if used on the
left-hand side of assignment, a field place. The \texttt{Expr::into\_place}
method converts suitable load, field, and array element expressions into a
\texttt{Place}. If conversion fails, the compiler reports ``expected a place
expression.''

This distinction is essential for assignment, compound assignment, address-of,
and pointer dereference. It also makes lower IR generation easier because places
can be lowered to pointers.

\section{Lowered Source Constructs}

Several source constructs are rewritten during \irtree{} lowering:

\begin{itemize}
  \item a \texttt{while} expression becomes a \texttt{loop} containing a
        conditional break;
  \item a \texttt{for} expression over a range becomes a block with hidden
        counter and target variables plus a normal \texttt{loop};
  \item \texttt{a || b} becomes a conditional expression that skips
        \texttt{b} when \texttt{a} is true;
  \item \texttt{a \&\& b} becomes a conditional expression that skips
        \texttt{b} when \texttt{a} is false;
  \item unary negation becomes subtraction from zero;
  \item function arguments become local variables initialized from argument
        expressions.
\end{itemize}

Lowering source constructs early reduces the number of forms later stages must
understand. Lower IR does not need a special \texttt{while} instruction, a
special \texttt{for} instruction, or a short-circuit boolean instruction. It only
needs blocks, loops, conditionals, arithmetic, comparisons, loads, stores, and
calls.

\section{Semantic Checkers}

After function bodies are lowered, the compiler runs semantic checkers. The
first checker verifies the \texttt{main} ABI. If a function named \texttt{main}
exists, it must have no arguments, must not be variadic, and must return
\texttt{i32}.

The second checker validates \texttt{comptime} expressions. It walks each
function body and checks the inner expression of each \texttt{comptime}. It
rejects operations that would depend on run-time state or unsafe effects.

The third checker validates \texttt{@pure} functions. A pure function must have
a body, may call only other pure functions, and may not use operations that are
not allowed in pure code. This check is needed because compile-time expressions
may call pure functions.

\section{Tree Optimization}

The \irtree{} optimizer is intentionally small and local. It performs
transformations that simplify generated trees without requiring global data-flow
analysis:

\begin{itemize}
  \item \texttt{\&ptr.*} becomes \texttt{ptr};
  \item \texttt{(\&place).*} becomes \texttt{place};
  \item empty blocks and single-expression blocks with no local variables are
        collapsed;
  \item code after the first expression of never type in a block is removed;
  \item pure unused expressions before the final expression are removed;
  \item redundant trailing unit constants are removed.
\end{itemize}

Listing \ref{lst:dead-elim-tree} shows an example. The source calls
\texttt{external()}, returns, and then calls \texttt{external()} again. The
second call is unreachable and is removed from the \irtree{} snapshot.

\begin{lstlisting}[caption={IR\_TREE after dead-code cleanup},label={lst:dead-elim-tree}]
fn test() -> unit
    BLOCK TYPE=unit
    |   FUNCTION_CALL("external") TYPE=unit
    |   RETURN TYPE=!
    |   |   CONST UNIT TYPE=unit
\end{lstlisting}

\section{IR\_TREE Example}

The \texttt{comptime\_primes.epl} example computes a table of twenty primes at
compile time. After compile-time evaluation, the tree no longer contains the
loop that computed the values. It contains a store of one constant array:

\begin{lstlisting}[caption={Compile-time result in IR\_TREE}]
STORE TYPE=unit
|   VARIABLE(var_1) [PLACE] TYPE=[i32; 20]
|   CONST [3, 5, 7, 11, 13, 17, 19, 23, 29, 31,
|          37, 41, 43, 47, 53, 59, 61, 67, 71, 73]
|          TYPE=[i32; 20]
\end{lstlisting}

This is the main reason \irtree{} is useful as both a semantic representation
and a testing artifact. It shows the result of lowering and compile-time
execution directly.
